% ======================================================================================================================
%   PAQUETES BÁSICOS
% ======================================================================================================================
\usepackage{polyglossia}       % Soporte multilenguaje (XeLaTeX/LuaLaTeX)
\setmainlanguage{spanish}      % Idioma principal
\usepackage{graphicx}          % Inserción de imágenes
\usepackage{amsmath, amssymb, amsfonts} % Símbolos matemáticos
\usepackage{float}             % Posicionamiento exacto de figuras/tablas
\usepackage{tikz}              % Gráficos vectoriales
\usepackage{eso-pic}           % Elementos en el fondo
\usepackage{changepage}        % Ajustar márgenes locales
\usepackage{enumitem}          % Listas personalizadas
% \usepackage{lipsum}          % Texto de prueba (eliminar en versión final)
\usepackage{makecell}
\usepackage{pgf}

% ======================================================================================================================
%   TIPOGRAFÍA Y FORMATO DE TEXTO
% ======================================================================================================================
\usepackage{fontspec}          % Fuentes del sistema (XeLaTeX/LuaLaTeX)
\setmainfont{Arial }  % Fuente principal

\usepackage{setspace}          % Interlineado
\setstretch{1.5}               % Interlineado 1.5

\setlength{\parindent}{0pt}    % Sin sangría en párrafos
\setlength{\parskip}{6pt}      % Espacio entre párrafos (opcional, se ve más limpio)

% ======================================================================================================================
%   MÁRGENES Y GEOMETRÍA
% ======================================================================================================================
\usepackage[
    left=3cm,
    right=2.54cm,
    top=2.54cm,
    bottom=2.54cm,
    footskip=1.5cm,
    showframe=false
]{geometry}

% ======================================================================================================================
%   ESTILOS DE CAPÍTULOS Y SECCIONES
% ======================================================================================================================
\usepackage{titlesec}

% --- Capítulos ---
\titleformat{\chapter}
{\bfseries\fontsize{14pt}{16pt}\selectfont}
{\thechapter}{1em}{}
\titlespacing*{\chapter}{0pt}{0pt}{0pt}

% --- Secciones ---
\titleformat{\section}
{\bfseries\fontsize{12pt}{14pt}\selectfont}
{\thesection}{1em}{}
\titlespacing*{\section}{0pt}{16pt}{6pt}

% --- Subsecciones ---
\titleformat{\subsection}
{\bfseries\fontsize{12pt}{14pt}\selectfont}
{\thesubsection}{1em}{}
\titlespacing*{\subsection}{0pt}{16pt}{6pt}


% ======================================================================================================================
%   ENCABEZADOS Y PIES DE PÁGINA
% ======================================================================================================================
\usepackage{fancyhdr}
\fancypagestyle{plain}{%
    \fancyhf{} % Limpia encabezados y pies
    \renewcommand{\headrulewidth}{0pt} % Sin línea superior
    \renewcommand{\footrulewidth}{0pt} % Sin línea inferior
    \cfoot{\thepage} % Número de página centrado
}
\pagestyle{plain} % Usar plain en todo el documento

\fancypagestyle{noheader}{%
    \fancyhf{} % Limpia encabezados y pies
    \renewcommand{\headrulewidth}{0pt}
    \renewcommand{\footrulewidth}{0pt}
    % Sin \cfoot{\thepage} ← esto quita el número
}

% ======================================================================================================================
%   HIPERVÍNCULOS
% ======================================================================================================================
\usepackage{hyperref}
\hypersetup{
    colorlinks=true,
    linkcolor=black,
    urlcolor=black,
    pdftitle={Tesina},
    pdfauthor={Tu Nombre},
}

% ======================================================================================================================
%   TABLA DE CONTENIDOS (con puntos)
% ======================================================================================================================
\usepackage{tocloft}
\renewcommand{\cftchapleader}{\cftdotfill{\cftdotsep}}
\renewcommand{\cftsecleader}{\cftdotfill{\cftdotsep}}
\renewcommand{\cftsubsecleader}{\cftdotfill{\cftdotsep}}


\addto\captionsspanish{ %Cambia el titulo del indice de cuadros
    \renewcommand{\listtablename}{Índice de tablas}
}

% ======================================================================================================================
%   NUEVAS MACROS
% ======================================================================================================================

% Capítulos sin número pero visibles en el índice
\newcommand{\chapterstar}[1]{%
  \chapter*{#1}%
  \addcontentsline{toc}{chapter}{#1}%
  \markboth{}{#1}%
}

% Página independiente para título de capítulo (sin número de página ni en el índice)
\newcommand{\chapterpage}[2]{%
  \clearpage
  \thispagestyle{empty} % Sin encabezado ni número
  \null\vfill
  \begin{center}
    {\Huge \textbf{Capítulo #1}} \\[3ex]
    {\LARGE #2}
  \end{center}
  \vfill\null
  \clearpage
  \markboth{}{#2} % Marca para encabezados
}

\newcommand{\chapterpageanex}[1]{%
  \clearpage
  \thispagestyle{empty} % Sin encabezado ni número
  \null\vfill
  \begin{center}
    {\Huge \textbf{#1}}
  \end{center}
  \vfill\null
  \clearpage
  \markboth{}{#1} % Marca para encabezados
}

\newcommand{\chapterpageavan}[1]{%
  \clearpage
  \thispagestyle{empty} % Sin encabezado ni número
  \null\vfill
  \begin{center}
    {\Huge \textbf{#1}}
  \end{center}
  \vfill\null
  \clearpage
  \markboth{}{#1} % Marca para encabezados
}

\newcommand{\frontmatterwithnonumber}{
    \pagestyle{empty}
    \setcounter{page}{0}
    \pagenumbering{gobble}
}

\newcommand{\mainmatterwithnumber}{
    \clearpage
    \pagestyle{plain}
    \setcounter{page}{1}
    \pagenumbering{arabic}
}

% Comando para restablecer completamente la numeración
\newcommand{\startmainmatter}{
    \clearpage
    \pagenumbering{arabic}
    \setcounter{page}{1}
    \pagestyle{plain}
    \addtocontents{toc}{\protect\cftpagenumberson{chapter}}
    \addtocontents{toc}{\protect\cftpagenumberson{section}}
    \addtocontents{toc}{\protect\cftpagenumberson{subsection}}
}

\usepackage{tikz}
\usetikzlibrary{arrows.meta, positioning, calc}


% Comando para imagenes
\usepackage{graphicx}
\usepackage{caption}
\usepackage{hyperref}
\usepackage{chngcntr} % Para numerar figuras por capítulo

% Figuras numeradas como Capítulo.Figura
\counterwithin{figure}{chapter}

% Texto "Fig." en lugar de "Figure"
\captionsetup[figure]{name=Fig.}

% Espacio entre marco e imagen
\setlength{\fboxsep}{6pt} % Puedes cambiar el valor
\setlength{\fboxrule}{0.5pt} % Grosor del marco

\usepackage{hyperref}
\usepackage{url}

% Tablas 
\usepackage{tabularx}
\usepackage{booktabs}
\usepackage{caption}
\usepackage{longtable}
\usepackage{array}
\usepackage[table]{xcolor}
\usepackage{xparse}  


\usepackage{xfp}   
\usepackage{pgf}    

\newcommand{\TotalGeneral}{0}
% Unidades en tablas
\usepackage{siunitx}  

% Configuración de siunitx para formato monetario
\sisetup{
    mode = text, 
    group-separator = {,},       % Separador de miles
    group-minimum-digits = 4,    % Agrupa a partir de 4 dígitos
    output-decimal-marker = {.}, % Marcador decimal
    round-mode = places,         % Redondea a un número fijo de decimales
    round-precision = 2,         % Usar 2 decimales
    round-integer-to-decimal = true, % Añade .00 a enteros (e.g., 80 -> 80.00)
}

\NewDocumentCommand{\fila}{m m m o}{%
    % #1: nombre del componente
    % #2: precio unitario (numérico)
    % #3: cantidad numérica (para el cálculo)
    % #4: (opcional) texto a mostrar en la columna cantidad
    \xdef\TotalGeneral{\fpeval{\TotalGeneral + (#2 * #3)}}%
    #1 & 
    \num{#2} & 
    \IfValueTF{#4}{#4}{\num{#3}} &   % Si hay opcional, muestra el texto; si no, muestra el número formateado
    \num{\fpeval{#2 * #3}} \\
    \hline
}



% imagenes
% Activar formato de figuras para avances
\newcommand{\avancesfigureson}{%
  \renewcommand{\figurename}{Av.}%
  \setcounter{figure}{0}%
  \renewcommand{\thefigure}{Av.\arabic{figure}}%
}

% Restaurar formato estándar de figuras (para capítulos numerados)
\newcommand{\avancesfiguresoff}{%
  \renewcommand{\figurename}{Fig.}%
  \counterwithin{figure}{chapter}%   Restablece la numeración por capítulo
  \setcounter{figure}{0}%
}


% Separador linea
\newcommand{\secbreak}{%
    \vspace{4pt}
    \noindent\rule{\textwidth}{0.4pt}
    \vspace{4pt}
}